\begin{recomendations}
% Revisión 2026-11-07: resultados actualizados y procedencia en Experiments.tex y fix.md.

La primera recomendación es \textbf{ampliar y diversificar la muestra de lectores}, incorporando
personas con otras formaciones y hábitos de lectura. Conviene conservar las respuestas por lector,
par y criterio, identificar al autor y comparar los resultados con y sin su participación. Alternar
el orden de presentación y la posición A/B permite estudiar los efectos de cansancio y posición.
Las nuevas comparaciones deberían conectar los cinco pares para poder estimar un ranking humano
global, además de preferencias dentro de cada pareja.

Se recomienda \textbf{completar la trazabilidad del estudio}: vincular los diez textos a sus
identificadores de ejecución, recuperar las respuestas individuales y documentar por separado las
lecturas adicionales de los mismos participantes. La comparación de memoria propia y compartida
comunicada por el autor debe acompañarse de los planes utilizados, los recuentos de menciones a
hechos desconocidos y las preferencias. Su replicación, junto con el contraste entre simulación y
escritura directa desde el mismo plan, permitirá estimar la estabilidad de las observaciones.

En el análisis horizontal se debe mantener la distinción entre ejecución completada y extracción
compatible. Es necesario \textbf{reunir las extracciones de las 24 unidades} bajo el mismo protocolo
y conservar las repeticiones que estiman el ruido. Para el extractor, una revisión humana de
historias completas debe registrar tanto las ocurrencias citadas como las omitidas y documentar
quién revisó cada caso. Así podrán contrastarse precisión, cobertura y repetibilidad sin confundir
la ausencia de una anotación con un resultado nulo.

La siguiente prioridad es \textbf{reproducir la validación del juez}: publicar sus aciertos y
denominadores por criterio y partición, conservar las predicciones y compararlo con la longitud
sola sobre los mismos votos. El rango de acierto comunicado debe vincularse con ese informe antes
de atribuirle generalización. Para ordenar el corpus de 122 historias se requieren su lista,
las extracciones compatibles, la identificación de los textos de entrenamiento y las razones de
exclusión; después corresponde analizar cualitativamente las primeras historias.

Como trabajo posterior se propone \textbf{usar el juez como función objetivo}: generar varias
historias con distintos parámetros y seleccionar la que el juez ordene primero, y después buscar
configuraciones mediante un \textbf{algoritmo genético}. Ambos usos requieren un juez que generalice
fuera de muestra y supere al control de longitud. La metodología también puede aplicarse a otros
textos generados, adaptando el catálogo de rasgos y las preguntas, sin sustituir las preferencias
humanas por un juicio de calidad del LLM.

\end{recomendations}
